\usepackage[utf8]{inputenc}
\usepackage[french]{babel}
\usepackage[T1]{fontenc}

\usepackage[margin=0.8in]{geometry}

\usepackage[svgnames]{xcolor}

\usepackage{amsmath, amssymb, amsfonts, amsthm}

\usepackage{tikz}
\usepackage{tkz-tab}

\usepackage{pgfplots}
\pgfplotsset{compat=1.9}


% theorems
\makeatother
\usepackage{thmtools}
\usepackage[framemethod=TikZ]{mdframed}
\mdfsetup{skipabove=1em,skipbelow=0em}


\theoremstyle{definition}

\declaretheoremstyle[
    headfont=\bfseries\sffamily\color{LimeGreen!70!black}, bodyfont=\normalfont,
    mdframed={
        linewidth=2pt,
        rightline=false, topline=false, bottomline=false,
        linecolor=LimeGreen, backgroundcolor=LimeGreen!5,
    }
]{thmgreenbox}

\declaretheoremstyle[
    headfont=\bfseries\sffamily\color{LimeGreen!70!black}, bodyfont=\normalfont,
    mdframed={
        linewidth=2pt,
        rightline=false, topline=false, bottomline=false,
        linecolor=LimeGreen
    }
]{thmgreenline}

\declaretheoremstyle[
    headfont=\bfseries\sffamily\color{DarkGreen!70!black}, bodyfont=\normalfont,
    mdframed={
        linewidth=2pt,
        rightline=false, topline=false, bottomline=false,
        linecolor=DarkGreen, backgroundcolor=DarkGreen!5,
    }
]{thmdarkgreenbox}

\declaretheoremstyle[
    headfont=\bfseries\sffamily\color{DarkGreen!70!black}, bodyfont=\normalfont,
    mdframed={
        linewidth=2pt,
        rightline=false, topline=false, bottomline=false,
        linecolor=DarkGreen
    }
]{thmdarkgreenline}

\declaretheoremstyle[
    headfont=\bfseries\sffamily\color{RoyalBlue!70!black}, bodyfont=\normalfont,
    mdframed={
        linewidth=2pt,
        rightline=false, topline=false, bottomline=false,
        linecolor=RoyalBlue, backgroundcolor=RoyalBlue!5,
    }
]{thmbluebox}

\declaretheoremstyle[
    headfont=\bfseries\sffamily\color{RoyalBlue!70!black}, bodyfont=\normalfont,
    mdframed={
        linewidth=2pt,
        rightline=false, topline=false, bottomline=false,
        linecolor=RoyalBlue
    }
]{thmblueline}

\declaretheoremstyle[
    headfont=\bfseries\sffamily\color{FireBrick!70!black}, bodyfont=\normalfont,
    mdframed={
        linewidth=2pt,
        rightline=false, topline=false, bottomline=false,
        linecolor=FireBrick, backgroundcolor=FireBrick!5,
    }
]{thmredbox}

\declaretheoremstyle[
    headfont=\bfseries\sffamily\color{FireBrick!70!black}, bodyfont=\normalfont,
    numbered=no,
    mdframed={
        linewidth=2pt,
        rightline=false, topline=false, bottomline=false,
        linecolor=FireBrick, backgroundcolor=FireBrick!1,
    },
    qed=QED
]{thmproofbox}

\declaretheoremstyle[
    headfont=\bfseries\sffamily\color{LimeGreen!70!black}, bodyfont=\normalfont,
    numbered=no,
    mdframed={
        linewidth=2pt,
        rightline=false, topline=false, bottomline=false,
        linecolor=LimeGreen, backgroundcolor=LimeGreen!1,
    }
]{thmsolutionbox}

\declaretheoremstyle[
    headfont=\bfseries\sffamily\color{RoyalBlue!70!black}, bodyfont=\normalfont,
    numbered=no,
    mdframed={
        linewidth=2pt,
        rightline=false, topline=false, bottomline=false,
        linecolor=RoyalBlue, backgroundcolor=RoyalBlue!1,
    },
]{thmexplanationbox}


\declaretheorem[style=thmbluebox, name=Définition]{definition}
\declaretheorem[style=thmgreenbox, numbered=no, name=Example]{example}
\declaretheorem[style=thmredbox, name=Proposition, numberlike=definition]{proposition}
\declaretheorem[style=thmredbox, name=Théoreme, numberlike=definition]{theorem}
\declaretheorem[style=thmredbox, name=Lemme, numberlike=definition]{lemma}

\declaretheorem[style=thmproofbox, name=Démonstration]{replacementproof}
\renewenvironment{proof}{\vspace{-12pt}\begin{replacementproof}}{\end{replacementproof}}

\declaretheorem[style=thmdarkgreenline, name=Solution, numbered=no]{solution_}
\newenvironment{solution}{\vspace{-12pt}\begin{solution_}}{\end{solution_}}

\declaretheorem[style=thmgreenline, numbered=no, name=Remarque]{remark}
\declaretheorem[style=thmgreenline, numbered=no, name=Note]{note}

\declaretheorem[style=thmbluebox, numbered=no, name=Notation]{notation}
\declaretheorem[style=thmgreenbox, numbered=no, name=Méthode]{method}
\declaretheorem[style=thmgreenbox, numbered=no, name=Rappel]{rappel}
\declaretheorem[style=thmdarkgreenbox, name=Exercice]{exercise}


\renewcommand{\emptyset}{\text{Ø}}

\newcommand{\C}{\ensuremath{\mathbb{C}}}
\newcommand{\R}{\ensuremath{\mathbb{R}}}
\newcommand{\Q}{\ensuremath{\mathbb{Q}}}
\newcommand{\D}{\ensuremath{\mathbb{D}}}
\newcommand{\Z}{\ensuremath{\mathbb{Z}}}
\newcommand{\N}{\ensuremath{\mathbb{N}}}
